\section{非torsion境界外圏と混合\texorpdfstring{$A_2$}{A2}の完全判定}
\label{sec:mixed-a2-category}
本節は，命題\ref{prop:actual-mixed-a2}のconnecting morphismを
導来完全忠実性を仮定せず決定する．さらに二つの根sheafが生成する
Serre部分圏全体を記述する．入力はRains
\cite[Corollaries 7.7, 7.13, 7.15, Propositions 7.6, 7.32,
Corollary 10.12, Proposition 10.13]{RainsBirational}，
前節までの符号別Ext，および境界triangleである．
$H=-K_X$はnef，$H^2=d>0$，$q$は非torsionとする．
構造line bundle $e$はexceptionalである．
以下，$\operatorname{Ext}_X$は常に\emph{ambientのRains heart}で計算する．

\subsection{境界外圏は有限長のSerre部分圏}
\begin{lemma}[零次元の境界外対象と商の排除]
\label{lem:offboundary-no-points}
$Li^*M=0$とする．$M$の零次元の商は零である．
従って，境界非交叉の既約純一次元sheafはRains heartの
\emph{単純対象}である．これは非torsion仮定を用いる結論であり，
既約性だけから全射性を推論するものではない．
\end{lemma}
\begin{proof}
零次元の商$M\twoheadrightarrow Z$を取る．
$i^*$は右完全なので$i^*Z=0$であり，境界triangleから
$\theta Z\to Z$は全射である．
Rainsの$\theta$の数値作用式によれば$[\theta Z]=[Z]$である．
同じ数値類の間の全射は同型
\cite[Corollary 7.7]{RainsBirational}だから$Li^*Z=0$となる．
一方，$[Z]=\ell[\mathrm{pt}]$，$\ell=\chi(e,Z)\ge0$であり，
その境界制限のdeterminantは$q^\ell$である
\cite[Proposition 7.6, \S7, Proposition 7.24の直前]{RainsBirational}．
非torsion性から$\ell=0$，従って$Z=0$である．
既約純一次元sheafの任意の非零部分対象の商は零次元なので，
後半も従う．
\end{proof}

\begin{proposition}[境界外圏の有限長性]
\label{prop:offboundary-serre-length}
$\mathcal Z=\{M\in\mathcal A_{\mathrm R}:\dim M\le1,\ Li^*M=0\}$は
Serre部分圏であり，全対象は有限長である．
その単純対象は境界非交叉の2-sphericalな根sheafである．
最大共鳴で全単純根sheafが実在する場合，その単純対象は
定理\ref{thm:only-simple-irreducible}の$r=9-d$個に尽きる．
\end{proposition}
\begin{proof}
$0\to N\to M\to P\to0$，$M\in\mathcal Z$とする．
次元は部分商で増加しない．$H$のnef性とChern類の加法性から
\[
0=Hc_1(M)=Hc_1(N)+Hc_1(P),\qquad Hc_1(N),Hc_1(P)\ge0,
\]
よって両項は零である．$i^*P=0$なので$\theta P\to P$は全射．
rank零かつ$K_Xc_1(P)=0$ゆえ，Rains Corollary 7.15により
$[\theta P]=[P]$であり，Corollary 7.7から同型である．
従って$Li^*P=0$．短完全列における$\theta\to\mathrm{id}$の自然性から
$\theta N\to N$も同型である．拡大閉性は境界triangleから従う．

非零零次元対象は補題\ref{lem:offboundary-no-points}により存在しない．
従って真の降鎖の非零逐次商はすべて一次元であり，
Rains Proposition 7.32のfinitely partitive性によって降鎖は停止する．
coherent heartのnoetherian性と合わせて有限長性を得る．

単純対象$S$について$D=c_1(S)\ne0$，$HD=0$である．
$H^\perp$は負定値で，Riemann--Rochの偶奇から$D^2$は偶数．
$\operatorname{End}(S)=k$，$\theta S=S$とSerre双対性から
\[
-D^2=\chi(S,S)=2-\dim\operatorname{Ext}_X^1(S,S)\le2.
\]
従って$D^2=-2$，$\operatorname{Ext}_X^1(S,S)=0$であり，根sheafとなる．
最後の主張には既に証明した単純根の排他性と一意性を使う．
\end{proof}

\subsection{Euler指数和が一の場合のconnecting map}
$S_+=M_\alpha$，$S_-=M_\beta$を境界非交叉の既約根sheafとし，
\[
\alpha\beta=1,\qquad m_+=p,\qquad m_-=-(p-1),\qquad p\ge2
\]
とする．$\gamma=\alpha+\beta$は根で，$m_\gamma=1$である．
両向きの$\operatorname{Ext}_X^1$は一次元なので，非分裂拡大を
\begin{equation}\label{eq:a2-two-extensions}
0\to S_-\to P\to S_+\to0,\qquad
0\to S_+\to Q\to S_-\to0
\end{equation}
と書く．各中間項は同型を除き一意である．
ここで$P,Q$はsheafの記号であり，境界は引き続き$E$と書く．

\begin{theorem}[connecting mapの無条件な最大rank]
\label{thm:a2-delta-rank}
非分裂拡大$P$のconnecting mapは全射であり，
\[
\delta:V_+\twoheadrightarrow V_-,\qquad
\operatorname{rank}\delta=p-1,\qquad
(\dim V^0(P),\dim V^1(P))=(1,0).
\]
さらにline bundle $L$を核とする短完全列
\[
0\longrightarrow L\longrightarrow e\longrightarrow P\longrightarrow0,
\qquad c_1(L)=-\gamma,\quad \chi(e,L)=0
\]
が存在する．特に命題\ref{prop:actual-mixed-a2}の$p=2$では
$\delta:k^2\to k$のrankは\emph{1}であり，
$V_+$は1次元加群$V_-$を商に持つ．この結論に
$\Phi$の完全忠実性や全$q$導来同値は必要ない．
\end{theorem}
\begin{proof}
符号別Extから$\dim\Hom(e,S_+)=p$，
$\dim\Ext^1(e,S_-)=p-1$，$\Hom(e,S_-)=0$である．
長完全列により$\Hom(e,P)=\ker\delta\ne0$．
非零射$f:e\to P$を取ると，$P\to S_+$との合成は非零なので全射．
その像$I$と$S_-$の交わりは，単純性により$0$か$S_-$である．
前者なら$I\cong S_+$が拡大を分裂させるため不可能であり，
後者なら$I=P$．従って$f$は全射である．

$L=\ker f$はrank-one torsion-freeで，
$c_1(L)=-\gamma$，$\chi(e,L)=1-\chi(e,P)=0$．
RainsのEuler式から
\[
\chi(L,L)=-1+2\chi(e,L)-c_1(L)(c_1(L)-K_X)
=-1-(-2)=1.
\]
従ってRains Proposition 10.13の等号条件により$L$はline bundle．
$\Hom(e,L)=0$である：非零射があれば$L\hookrightarrow e$との合成は
$e$の非零自己射，従って同型となり$P\ne0$に反する．
また
\[
\Ext^2(e,L)=\Hom(L,\theta e)^*=0.
\]
実際，右辺の非零射があれば
$c_1(\theta e)-c_1(L)=K_X+\gamma$が有効となる
\cite[Corollary 10.12]{RainsBirational}が，
$H(K_X+\gamma)=-d<0$はnef性に反する．
$\chi(e,L)=0$なので$\Ext^1(e,L)=0$も従う．
$e$のexceptionalityと$0\to L\to e\to P\to0$から
$\Hom(e,P)=k$，$\Ext^1(e,P)=0$を得る．
従って$\delta$は全射である．
\end{proof}

\begin{corollary}[直和を用いた二段拡大の正確な公式]
\label{cor:a2-rank-matrix-exact}
$0\to S_-^{\oplus b}\to M\to S_+^{\oplus a}\to0$の拡大類を
$T\in\operatorname{Mat}_{b\times a}(k)$とし，$t=\operatorname{rank}T$とする．
すると
\[
\operatorname{rank}\delta_M=(p-1)t,\qquad
\dim V^0(M)=pa-(p-1)t,\qquad
\dim V^1(M)=(p-1)(b-t).
\]
$p=2$では$(\dim V^0,\dim V^1)=(2a-t,b-t)$である．
\end{corollary}
\begin{proof}
$\delta_M=\delta\otimes T$と前定理を使う．
\end{proof}

\subsection{混合\texorpdfstring{$A_2$}{A2}の全有限長対象とtorsion pair}
\begin{theorem}[四つの直既約対象による完全分類]
\label{thm:a2-serre-category}
$\mathcal C=\langle S_+,S_-\rangle_{\rm Serre}$をRains heartの
有限長Serre部分圏とする．$\mathcal C$の直既約対象は
\[
S_+,\quad S_-,\quad P,\quad Q
\]
の四つだけである．$P\oplus Q$は$\mathcal C$のprojective generatorであり，
\[
\mathcal C\simeq\operatorname{mod}_{fd}\Pi(A_2),\qquad
\Pi(A_2)=k(\, +\mathrel{\substack{\longrightarrow\\[-2pt]\longleftarrow}}-\,)/
(\hbox{両方の長さ二の道})
\]
となる．これは次元4のradical-square-zero代数である．
\end{theorem}
\begin{proof}
非分裂性と単純性から，順序$(S_+,S_-,P,Q)$でのHom次元表は
\[
\left(\dim\Hom(U,V)\right)_{U,V}=
\begin{pmatrix}1&0&0&1\\0&1&1&0\\1&0&1&1\\0&1&1&1\end{pmatrix}.
\]
例えば$\Hom(S_+,P)=\Hom(P,S_-)=0$は，
$k\to\operatorname{Ext}_X^1(S_+,S_-)=k$のconnecting mapが
拡大類を送る同型であることから従う．
$\Hom(P,Q)=k$は$P\twoheadrightarrow S_+\hookrightarrow Q$で生成され，
逆方向も同様で，両合成は零である．

全対象で$\theta$が自然に恒等なので，Serre双対性とEuler形式により
\[
\dim\operatorname{Ext}_X^1(U,V)
=\dim\Hom(U,V)+\dim\Hom(V,U)+c_1(U)c_1(V).
\]
この表を代入すると，$S_+,S_-$間の両向きの一次元Ext以外は全て零である．
とくに$\operatorname{Ext}_X^1(P,S_\pm)=
\operatorname{Ext}_X^1(Q,S_\pm)=0$．
長さに関する帰納法から$P,Q$は$\mathcal C$でprojectiveであり，
各単純対象への全射を持つので$P\oplus Q$がprojective generatorとなる．
通常のMorita理論により，$\mathcal C$はそのEnd代数の有限次元右加群圏と同値．
上記のHomと合成の記述が次元4の表示を与える．

最後にこの代数の表現を，$U_+,U_-$と二写像
$x:U_+\to U_-$，$y:U_-\to U_+$，$xy=yx=0$で書く．
$\operatorname{im}y\subset\ker x$，$\operatorname{im}x\subset\ker y$を用い
各空間に適合する基底を選ぶと，$x$のみが同型の一次元対，
$y$のみが同型の一次元対，および孤立した一次元空間に分解できる．
これらがそれぞれ$P,Q,S_+,S_-$である．
\end{proof}

\begin{corollary}[局所torsion pairと全対象での非消失]
\label{cor:a2-torsion-pair}
$\mathcal C$では
\[
\mathcal T=\ker\Ext^1(e,-)=\operatorname{add}(S_+,P),\qquad
\mathcal F=\ker\Hom(e,-)=\operatorname{add}(S_-)
\]
がtorsion pairである．各直既約対象のcohomology次元は
\[
\begin{array}{c|rrrr}
M&S_+&S_-&P&Q\\\hline
\dim V^0(M)&p&0&1&p\\
\dim V^1(M)&0&p-1&0&p-1
\end{array}
\]
となる．特に$\Phi(M)=0$である$M\in\mathcal C$は$M=0$だけである．
一般の上記quiver表現では$a=\dim U_+$，$b=\dim U_-$，
$t=\operatorname{rank}x$として
\[
\dim V^0(M)=pa-(p-1)t,\qquad
\dim V^1(M)=(p-1)(b-t).
\]
$p=2$では$\chi(e,P)=\chi(e,Q)=1$だが，
$P$の像は次数0に集中し，$Q$の像は二つの次数に残る．
従って非既約対象のEuler指数の符号だけではheartの位置を判定できない．
\end{corollary}
\begin{proof}
$P$の項は定理\ref{thm:a2-delta-rank}，$Q$の項は
第二の短完全列\eqref{eq:a2-two-extensions}の六項完全列から得る．
Hom表から$\Hom(\mathcal T,\mathcal F)=0$である．
$Q$については$0\to S_+\to Q\to S_-\to0$がtorsion分解を与え，
残る三対象は既にどちらかの部分圏に入る．直和を取れば全対象の分解を得る．
quiverの分解の重複度はそれぞれ
$a-\operatorname{rank}x-\operatorname{rank}y$，
$b-\operatorname{rank}x-\operatorname{rank}y$，
$\operatorname{rank}x$，$\operatorname{rank}y$であり，表から公式が従う．
非零直既約はいずれも像が非零なので最後の非消失も従う．
\end{proof}

\begin{remark}[局所分類と導来比較の区別]
このtorsion pairは$\mathcal C$上の結果であり，全$\mathcal Z$や全Rains heartへ
拡張したとは主張しない．また$\mathcal C$の\emph{内在的}高次ExtとambientのExtを
同一視してはならない．実際，$\Pi(A_2)$の単純加群のprojective resolutionは
二周期的で，$\operatorname{Ext}_{\mathcal C}^{2n}(S_+,S_+)=k$ $(n\ge0)$である．
一方$\operatorname{Ext}_X^j=0$ $(j>2)$なので，
$D^b(\mathcal C)\to D^b(\mathcal A_{\mathrm R})$は完全忠実でない．
本節は$\Phi$が全導来圏で同値であることも，$\mathcal C$上で射をすべて保つことも
証明していない．
\end{remark}

\subsection{残る二次元アフィン表現テストの縮約}
\begin{proposition}[完全忠実性に残る一つの代数的退化]
\label{prop:a2-scalar-test}
$p=2$とし，$W=V^0(P)$とする．有限次元右$A$加群の短完全列
\[
0\to W\to V_+\xrightarrow{\delta}V_-\to0,
\qquad \dim W=\dim V_-=1
\]
が存在し，$\delta$は非零である．さらに
\[
\dim\Hom_A(V_+,V_-)=
\begin{cases}2,&V_+\cong V_-^{\oplus2},\\1,&\text{それ以外}\end{cases}
\]
となる．従って自然写像
$\operatorname{Ext}_X^1(S_+,S_-)\to\Hom_A(V_+,V_-)$は常に単射であり，
同型であることは$V_+\not\cong V_-^{\oplus2}$と同値である．
\end{proposition}
\begin{proof}
短完全列は既知の六項完全列と$V^1(P)=0$による．
$\Hom_A(V_+,V_-)$は二次元の$\Hom_k(V_+,V_-)$の非零部分空間．
独立な二つの準同型があれば，それらを並べた写像
$V_+\to V_-^{\oplus2}$はベクトル空間としても$A$加群としても同型である．
逆は明らか．拡大類の生成元は非零$\delta$に写るので単射性も従う．
\end{proof}

右加群を行ベクトルで表し，$W$を第一座標，$\delta$を第二座標への射影に
取れば，作用は必ず
\[
\rho(a)=\begin{pmatrix}\chi_W(a)&0\\ \eta(a)&\chi_-(a)\end{pmatrix},\qquad
\eta(ab)=\eta(a)\chi_W(b)+\chi_-(a)\eta(b)
\]
の形になる．上の未解決退化は$\chi_W=\chi_-$かつ$\eta=0$に正確に対応する．
これは表現の\emph{正規形と検査条件}であって，指定楕円データに対する
$A$の生成元の$\chi_W,\chi_-,\eta$を算出したという主張ではない．
それらの計算と，全$q$導来比較の監査が次の課題である．

\begin{proposition}[characterの分離と非分裂性という必要条件]
\label{prop:a2-character-obstructions}
$p=2$で，$\Phi$が$S_+,S_-$の生成する三角圏上で完全忠実であるなら，
\[
W\not\cong V_-,\qquad
0\to W\to V_+\to V_-\to0\text{ は非分裂}
\]
でなければならない．従って上の二次元表現での$A$の像は
全ての下三角$2\times2$行列からなる次元3の代数である．
これは必要条件であって，導来完全忠実性の十分条件とはしない．
\end{proposition}
\begin{proof}
Hom表とEuler式から$\Ext_X^1(P,S_-)=0$である．
完全忠実性があれば，$\Phi(P)=W$，$\Phi(S_-[1])=V_-$なので
$\Hom_A(W,V_-)=0$，従って二つのcharacterは異なる．
また非零の包含$S_-\hookrightarrow P$は，$\Phi$の下で
$\Ext_A^1(V_-,W)$の元，すなわち表示した短完全列の拡大類へ移る．
忠実性からその類は非零である．

characterが異なると，像代数には対角成分の異なる行列があり，
その多項式から二つの相補的冪等元を得る．基底をそれらに合わせると，
非分裂性は像代数の非対角成分が非零であることに等しい．
この成分も一次元なので，二つの冪等元と合わせて全下三角代数を得る．
\end{proof}
\begin{remark}
$\chi_W=\chi_-$だが$\eta\ne0$というJordan型なら
命題\ref{prop:a2-scalar-test}の一次元Homテストには合格しても，
本命題のcharacter分離に失敗する．またcharacterが異なっても拡大が分裂すれば，
非零sheaf射が$\Phi$で零になる．対象の非消失と射の忠実性は別の条件である．
\end{remark}
